\chapter*{핵심 정리}
\addcontentsline{toc}{chapter}{핵심 정리}
\markboth{핵심 정리}{}
\begingroup
\captionsetup{labelformat=empty}

\begin{keybox}[한 문장 결론]
\textbf{LendingClub 대출은 2013--16년에 전부 사도 국채보다 나았다. 그래서 질문은 ``무엇을 고르느냐''다. 우리는 대출마다 부도확률과 금리로 기대수익을 계산해, 승인비율이 몇 \%든 기대수익이 높은 대출부터 사자고 제안한다. train에서 승인비율 10--90\% 9개 모두 같은 양을 무작위로 산 것보다 더 벌었고, 승인비율 60\% 이상이면 덜 흔들리기도 했다.}
\end{keybox}
숫자는 모두 train(2013--16년 발행 36개월 대출, 48개 발행월) 결과이며, test는 아직 열지 않았다. 판정 기준은 test 전에 고정했다(규칙서 v2.3).

\begin{formulabox}[승인비율 $q$란]
매달 그 달에 나온 36개월 대출을 점수 순으로 줄 세워 위에서부터 사는 \textbf{건수의 비율}이다(10\%, 20\%, \dots, 90\%). 투자자가 그 달 시장의 몇 \%를 살 돈이 있는지를 나타내며, 사지 않은 돈은 국채에 둔다고 본다. 대출마다 금액이 달라 \textbf{금액 비중} $w$는 $q$와 조금 다르다(예: 부도확률 순 80\%에서 평균 82.5\%). 선별 가치는 승인한 1달러당 IRR 차이로, Sharpe는 금액 비중 $w$를 곱한 초과수익으로 계산한다.
\end{formulabox}

\begin{formulabox}[선별 가치와 포트폴리오 초과수익은 다르다]
\textbf{선별 가치}는 \textbf{대출에 넣은 돈의 질}을 잰다. 어느 달 대출이 1,000건이고 승인비율이 10\%면, 점수 상위 100건과 아무 100건을 각각 사서 대출이 끝난 뒤 두 묶음의 실제 IRR을 비교한다. 예를 들어 우리 묶음 연 8.0\%, 무작위 묶음 연 5.9\%면 선별 가치는 2.1\%p다. IRR은 ``넣은 1달러가 1년에 몇 \% 벌었나''이므로, 얼마를 넣었는지와 상관없이 \textbf{같은 양을 무작위로 산 것보다 1달러당 몇 \%p 더 버나}를 뜻한다. 48개월 평균이 표의 숫자다. 고르는 실력을 재는 숫자이고, 주 판정에 쓴다.

\textbf{포트폴리오 초과수익}은 \textbf{전체 돈의 성과}다. 사지 않은 돈은 국채에 두어 초과수익이 0이므로, 승인비율 10\%면 돈의 90\%가 국채보다 더 벌지 못한다. 투자자가 실제로 국채보다 얼마나 더 버는지는 이 숫자가 보여 준다.

\smallskip
\centering\small
\begin{tabular}{@{}lccccc@{}}\toprule
기대수익 순(V1), train & 10\% & 50\% & 80\% & 90\% & 전부 승인\\\midrule
선별 가치(1달러당, 무작위 대비) & +2.11 & +0.99 & +0.48 & +0.28 & 0\\
포트폴리오 초과수익(국채 대비, 전체 돈) & 0.65 & 2.98 & 4.34 & 4.67 & 4.86\\\bottomrule
\end{tabular}

\raggedright
\smallskip
\textbf{읽는 법:} 적게 살수록 가장 좋은 대출만 골라 1달러당 차이가 크지만, 돈 대부분이 국채에 있어 전체 수익은 작다. 많이 살수록 평범한 대출까지 사게 되어 1달러당 차이는 줄지만 전체 수익은 커진다. 그래서 ``10\%만 사는 것이 가장 좋다''는 뜻이 아니다. 이 기간에는 대출이 국채보다 나았으므로 많이 살수록 전체 수익은 컸고, 우리 주장은 \textbf{몇 \%를 사든 같은 양을 무작위로 사는 것보다 낫다}는 것이다.
\end{formulabox}


\section*{흐름}
\begin{figure}[H]\centering
\begin{tikzpicture}[x=1mm,y=1mm,box/.style={draw=black!70,rounded corners=1.2mm,align=center,font=\footnotesize,minimum height=12mm,text width=32mm,fill=white,execute at begin node={\XeTeXlinebreakskip=0pt\relax}},ar/.style={-{Stealth},black!65}]
\node[box] (a) at (0,0) {그 달 새 대출\\(36개월 만기)};
\node[box] (b) at (42,0) {부도확률 $p$\\시점 학습 LightGBM\\(그날 알던 정보만)};
\node[box] (c) at (84,0) {기대수익 점수\\$s=(1-p)r-p\cdot\mathrm{LGD}$\\$-$ 국채 금리};
\node[box,line width=0.9pt] (d) at (126,0) {$s$ 높은 순으로\\예산만큼 구매\\남는 돈은 국채};
\draw[ar] (a)--(b); \draw[ar] (b)--(c); \draw[ar] (c)--(d);
\draw[dashed,black!35] (-18,-13)--(144,-13);
\node[font=\scriptsize\itshape,black!60,anchor=south west] at (-18,-13) {위: 승인하는 날(결과를 모름)};
\node[font=\scriptsize\itshape,black!60,anchor=north west] at (-18,-13) {아래: 대출이 끝난 뒤(실제 결과로 채점, $s$는 쓰지 않음)};
\node[box] (e) at (126,-27) {우리가 산 묶음\\실제 현금흐름\\\ensuremath{\to} 코호트 IRR};
\node[box] (f) at (126,-47) {같은 양 무작위(10회)\\실제 현금흐름\\\ensuremath{\to} 코호트 IRR};
\node[box,line width=0.9pt] (g) at (84,-37) {선별 가치\\$=$ 우리 IRR $-$ 무작위 IRR\\(국채는 양쪽에서 지워짐)};
\node[box] (h) at (42,-37) {48개 발행월 평균\\예산 10--90\%\\블록 부트스트랩 구간};
\node[box,line width=0.9pt] (i) at (0,-27) {주 판정(H2)\\95\% 구간 하한 $>0$\\\ensuremath{\to} 제안 유지};
\node[box,dashed] (j) at (0,-47) {보조 지표\\포트폴리오 Sharpe\\(평균 \ensuremath{\div} 흔들림)};
\draw[ar] (d)--(e); \draw[ar] (e)--(g); \draw[ar] (f)--(g); \draw[ar] (g)--(h); \draw[ar] (h)--(i); \draw[ar,dashed] (h)--(j);
\end{tikzpicture}
\caption{제안의 흐름. 위 줄은 승인하는 날 하는 일, 아래 줄은 대출이 끝난 뒤 채점하는 일이다. 코호트 IRR은 빌려준 돈(0개월), 월 납입금, 마지막 달의 조기상환 일시금이나 부도 후 회수금을 묶음 전체로 더해 구하고, 받은 돈에서 수수료 1\%를 뗀다(2장).}
\end{figure}


\section*{핵심 기술과 근거}
\begin{table}[H]\centering\small
\begin{tabularx}{\linewidth}{@{}>{\raggedright\arraybackslash}p{30mm}>{\raggedright\arraybackslash}p{50mm}L@{}}\toprule
기술 & 무엇을, 왜 & 근거(train 결과)\\\midrule
시점 학습(as-of) & 연도 Y 대출의 $p$는 Y년 1월 1일까지 결과가 확정된 대출로만 학습한 모델로 계산. 그때 투자자가 가질 수 있던 정보만 쓰기 위해 & 무작위 분할 AUC 0.73 vs 시점 학습 0.62--0.68. 차이가 ``미래를 본 값''(3장)\\
LightGBM 부도확률 $p$ & 작은 나무를 차례로 쌓아 앞 나무가 틀린 부분을 고침. 결측을 그대로 받고 변수 조합을 자동으로 잡음. 101개 특성, 시드 3개 평균 & 17종 비교에서 XGBoost와 선별 가치 같음(1.03 vs 1.04\%p). AUC 1위 랜덤포레스트는 돈을 덜 범(0.85\%p). 성과가 특정 모델에 달려 있지 않음(7장)\\
$p$는 순위로 사용 & $p$의 크기는 실제보다 낮게 나오지만 순서는 유지 & 2013년 예측 6.1\% vs 실제 12.2\%. 실제 수준으로 맞춰 다시 골라도 승인 대상 88--92\% 같음(5장)\\
기대수익 점수 $s$ & $s=(1-p)r-p\cdot\mathrm{LGD}-$국채. 예상 손실을 빼고도 국채보다 많이 남는 대출을 찾기 위해(9장 ``왜 $p$가 아니라 $s$인가) & 같은 등급 안 $s$ 상위 $-$ 하위 연 +2.38\%p(2.08--2.73), 48개월 모두 양수(7장)\\
발행월 컷 & 매달 그 달 대출끼리 줄 세워 위에서부터 승인비율 $q$만큼 승인 & 시장 상황이 같은 대출끼리 비교\\
코호트 IRR & 같은 달에 산 대출 묶음의 실제 현금흐름(투자, 월 납입금, 조기상환·\allowbreak{}회수금)으로 구한 연수익률, 수수료 1\% & 단순 연환산이나 끝난 대출만 보면 결론이 뒤집힘(2장)\\
선별 가치(주 판정) & 같은 달, 같은 양을 무작위(10회)로 산 묶음과의 IRR 차이. 국채는 양쪽에서 지워짐 & 아래 결과 표\\
포트폴리오 Sharpe(보조) & 매달 (IRR $-$ 국채) \ensuremath{\times} 승인 비중의 평균 \ensuremath{\div} 표준편차. 덜 흔들리는가 & 아래 결과 표. 약 3으로 큰 것은 만기 보유와 좋은 시기 때문이라 전략끼리 비교에만 사용\\
블록 부트스트랩 & 48개월을 6개월 덩어리로 2000번 다시 뽑아 구간 계산. 이웃 달의 상관을 보존 & 모든 전략 공통 인덱스, 짝지은 차이\\
정규화 & LightGBM에는 불필요(순서로만 나눔). 비교용 로지스틱에는 학습 데이터 기준 자르기, 중앙값 채우기, 표준화 & 표준화 폭발(최대 290, AUC 0.508)을 가용성 규칙으로 해결(0.678)\\
사전 고정과 정정 기록 & test 판정 기준을 test 전에 문서로 고정하고 git 태그. 틀린 설명은 정정표에 남김 & 규칙서 v1--v2.3, 정정표(8장)\\\bottomrule
\end{tabularx}
\caption{핵심 기술. 괄호의 장은 본문 위치.}
\end{table}

\section*{결과: 승인비율별 선별 가치와 Sharpe}
\begin{table}[H]\centering\small
\begin{tabular}{@{}rccccccc@{}}\toprule
 & \multicolumn{2}{c}{선별 가치(\%p)} & \multicolumn{3}{c}{포트폴리오 Sharpe} & \multicolumn{2}{c}{V1 $-$ 무작위}\\\cmidrule(lr){2-3}\cmidrule(lr){4-6}\cmidrule(l){7-8}
승인비율 & \textbf{V1(제안)} & V2 & \textbf{V1} & V2 & 무작위 & 선별 가치 90\% & Sharpe 차이 90\%\\\midrule
10\% & \textbf{+2.11} & $-0.32$ & 2.55 & 3.82 & 2.63 & 1.80, 2.45 & $-0.59$, 0.14\\
20\% & \textbf{+1.74} & $-0.04$ & 2.75 & 3.61 & 2.68 & 1.52, 1.99 & $-0.47$, 0.40\\
30\% & \textbf{+1.46} & +0.18 & 2.91 & 3.34 & 2.68 & 1.29, 1.62 & $-0.28$, 0.58\\
40\% & \textbf{+1.19} & +0.30 & 3.01 & 3.17 & 2.70 & 1.06, 1.33 & $-0.12$, 0.61\\
50\% & \textbf{+0.99} & +0.37 & 3.00 & 3.12 & 2.72 & 0.84, 1.12 & $-0.04$, 0.55\\
60\% & \textbf{+0.82} & +0.39 & 3.04 & 3.11 & 2.72 & 0.68, 0.94 & \textbf{0.08, 0.54}\\
70\% & \textbf{+0.63} & +0.37 & 3.05 & 3.07 & 2.73 & 0.51, 0.74 & \textbf{0.15, 0.52}\\
80\% & \textbf{+0.48} & +0.32 & 3.01 & 3.06 & 2.72 & 0.39, 0.56 & \textbf{0.15, 0.46}\\
90\% & \textbf{+0.28} & +0.21 & 2.93 & 2.98 & 2.72 & 0.22, 0.33 & \textbf{0.14, 0.36}\\\midrule
평균 & \textbf{+1.08} & +0.20 & & & & & \\\bottomrule
\end{tabular}
\caption{V1 = 기대수익 순(제안), V2 = 부도확률 낮은 순(대안, 이전 최종안). 발행월 컷, 국채 기준, 수수료 1\%, 무작위 10회 평균, 6개월 블록 부트스트랩 2000회(42단계, \ft{results/42_final_proposal_check.csv}). 선별 가치 구간은 모두 0보다 위, Sharpe 차이의 굵은 구간도 0보다 위.}
\end{table}

\begin{enumerate}[leftmargin=*]
\item \textbf{더 벌기(주 판정):} V1은 9개 승인비율 모두 무작위보다 더 벌었다(구간 모두 0보다 위). 승인한 1달러당 이점(선별 가치)은 승인비율이 작을수록 크다(10\%: +2.11, 80\%: +0.48\%p). 다만 포트폴리오 전체 초과수익은 많이 살수록 크다(10\%: 0.65, 80\%: 4.34\%p).
\item \textbf{덜 흔들리기(보조):} V1의 Sharpe는 승인비율 60\% 이상에서만 무작위보다 확실히 높다. 50\% 이하에서는 구분되지 않고, 10\%는 점추정이 조금 낮다.
\item \textbf{맞교환:} V2는 Sharpe가 모든 승인비율에서 가장 높지만 더 버는 효과는 작다(평균 +0.20\%p). 덜 흔들리는 것이 가장 중요한 투자자에게는 V2를 대안으로 보고한다. 승인비율이 클 때(70--90\%)는 두 전략의 결과가 거의 같다.
\item \textbf{기준 수익률:} 회사채(BAA) 기준에서는 Sharpe도 V1이 높다(승인 30\%: 1.29 vs V2 0.92, 무작위 0.52).
\end{enumerate}

\section*{test에서 확인하는 것}
\begin{itemize}[leftmargin=*]
\item \textbf{주 판정(v1의 H2 그대로):} 기대수익 순위의 선별 곡선(승인비율 10--90\%)이 무작위보다 위에 있는 면적의 95\% 구간 하한이 0보다 큰가. 통과하면 제안 유지, 실패하면 철회.
\item \textbf{보고만:} 위 표의 승인비율별 숫자, V2 80\%의 무작위 대비 Sharpe(이전 판정 H10), 기준 수익률 4단계. 결과가 어떻든 제안을 바꾸지 않는다.
\end{itemize}

\section*{핵심 한계}
\begin{itemize}[leftmargin=*]
\item 모든 숫자는 train이라 낙관적일 수 있다.
\item 승인비율이 클 때는 1달러당 이점이 작고(80\%: +0.48\%p), 승인비율이 작을 때 덜 흔들린다는 근거는 없다.
\item 2013--16년에는 경기침체가 없었다. 불경기에 모델의 위험 순서가 유지되는지는 확인하지 못했다(부도끼리의 상관, 10장).
\item 주 판정을 v1의 H2로 되돌린 것은 train 결과를 본 뒤의 결정이다. 근거는 과제 공지 재확인(Sharpe 필수 아님)이며, 판정 정의는 바꾸지 않았다(9장).
\end{itemize}

\section*{이름표}
\begin{table}[H]\centering\small
\begin{tabularx}{\linewidth}{@{}p{26mm}L@{}}\toprule
이름 & 뜻\\\midrule
V1 & 기대수익 순. \textbf{최종 제안}\\
V2 & 부도확률 낮은 순. 덜 흔들리기 대안(이전 최종안: V2 80\%)\\
V3 & 같은 등급 안에서 최근 24개월 개설 계좌가 적은 순. train을 보고 고른 도전자\\
v1--v2.3 & test 판정 규칙서의 버전. 현재 유효한 것은 v2.3(주 판정 H2)\\
선별 가치 & 같은 달, 같은 양 무작위 대비 IRR 차이(승인 1달러당, 연 \%p)\\
승인비율 $q$ & 그 달 대출 중 사는 건수의 비율. 투자자가 가진 돈에 달려 있고 우리는 정하지 않음. 금액 비중 $w$와 조금 다름\\\bottomrule
\end{tabularx}
\end{table}
\endgroup
